\subsection{Diagonal endomorphisms}

DEFINITION 1. --- Let \(\mathrm{B} = (e_i)_{i \in \mathrm{I}}\) be an orthonormal basis of E. An endomorphism \(u\) is said to be diagonal in the basis B or diagonal with respect to B if there exists a family \(\lambda = (\lambda_i)_{i \in \mathrm{I}}\) of elements of K such that \(u(e_i) = \lambda_i e_i\) for all \(i \in \mathrm{I}\).

Let \(\mathrm{B} = (e_i)_{i \in \mathrm{I}}\) be an orthonormal basis of E. We denote by \(\mathcal{D}_{\mathrm{B}}(\mathrm{E})\) the set of endomorphisms of E that are diagonal with respect to B. This is a closed, unital commutative subalgebra of \(\mathcal{L}(\mathrm{E})\) . Let \(u \in \mathcal{D}_{\mathrm{B}}(\mathrm{E})\) . The family \(\lambda = (\lambda_i)_{i \in \mathrm{I}}\) such that \(u(e_i) = \lambda_i e_i\) is determined uniquely by u and by the basis B, and it is called the family of eigenvalues of u relative to B.

Suppose that K = R and identify E with a subspace of \(\mathrm{E}_{(\mathbf{C})}\) . Let \(\mathrm{B} = (e_i)_{i \in \mathrm{I}}\) be an orthonormal basis of E. Denote by \(\mathrm{B}_{(\mathbf{C})}\) the orthonormal basis \((1 \otimes e_i)_{i \in \mathrm{I}}\) of \(\mathrm{E}_{(\mathbf{C})}\) (cf. EVT, V, p. 29, cor. 1). The map \(u \mapsto u_{(\mathbf{C})}\) induces an injective morphism of algebras from \(\mathcal{D}_{\mathrm{B}}(\mathrm{E})\) into \(\mathcal{D}_{\mathrm{B}_{(\mathbf{C})}}(\mathrm{E}_{(\mathbf{C})})\) ; its image is the set of \(u \in \mathcal{D}_{\mathrm{B}_{(\mathbf{C})}}(\mathrm{E}_{(\mathbf{C})})\) such that \(u(\mathrm{E}) \subset \mathrm{E}\) ; in other words, the set of diagonal endomorphisms u in the basis \(\mathrm{B}_{(\mathbf{C})}\) whose eigenvalues are real.

For every \(i \in \mathrm{I}\) , we denote by \(p_i\) the orthogonal projector of E with image \(\mathrm{Ke}_i\) . We have \(\| p_i \| = 1\) for all \(i \in \mathrm{I}\) . The endomorphism \(p_i\) is diagonal in the basis B, and the family of its eigenvalues is the family \((\delta_{ij})_{j\in\mathrm{I}}\) (Kronecker symbol, cf. A, II, p. 24).

Let \(\mathrm{B}^{\prime}=(f_{i})_{i\in\mathrm{I}}\) be an orthonormal basis of E and \(u:E\to E\) the isometric isomorphism such that \(u(f_{i})=e_{i}\) . Then \(\mathcal{D}_{\mathrm{B}^{\prime}}(\mathrm{E})=u^{-1}\mathcal{D}_{\mathrm{B}}(\mathrm{E})u\) .

We endow I with the discrete topology and denote by \(\mathcal{B}(\mathrm{I}) = \mathcal{C}_{b}(\mathrm{I}; \mathrm{K})\) the unital Banach algebra of bounded functions on I with values in K (example 2 of I, p. 17); if K = C, it is a star algebra (example 2 of I, p. 102).

Lemma 1. --- Let \(u \in \mathcal{D}_{\mathrm{B}}(\mathrm{E})\) and let \(\lambda = (\lambda_i)_{i \in \mathrm{I}}\) be the family of its eigenvalues. The family \(\lambda\) is bounded and one has \(\sup_i |\lambda_i| = \| u\|\) .

One has \(|\lambda_i| \leqslant \|u\|\) for all \(i\) , whence \(\sup_{i \in I} |\lambda_i| \leqslant \|u\|\) . Since moreover

\[
\| u (x) \| ^ {2} = \sum_ {i \in \mathrm{I}} | \lambda_ {i} | ^ {2} | \langle e _ {i} | x \rangle | ^ {2} \leqslant \left(\sup _ {i \in \mathrm{I}} | \lambda_ {i} | ^ {2}\right) \| x \| ^ {2}
\]

for all \(x \in \mathrm{E}\) (EVT, V, p. 22, prop. 5), it follows that \(\| u\| = \sup |\lambda_i|\) .

PROPOSITION 1. --- a) The map \(\alpha\) from \(\mathcal{D}_{\mathrm{B}}(\mathrm{E})\) into \(\mathcal{B}(\mathrm{I})\) that associates with a diagonal endomorphism \(u\) the family of eigenvalues of \(u\) is an isometric isomorphism of Banach algebras over K. If K = C, it is a morphism of involutive algebras.

\begin{enumerate}
\def\labelenumi{\alph{enumi})}
\setcounter{enumi}{1}
\item
  For any bounded family \(\lambda = (\lambda_i)_{i\in \mathrm{I}}\) the family \((\lambda_ip_i)_{i\in \mathrm{I}}\) is summable in the space \(\mathcal{L}(\mathrm{E})\) endowed with the topology of pointwise convergence, and the sum of this family is the unique mapping \(u\in \mathcal{D}_{\mathrm{B}}(\mathrm{E})\) such that \(\alpha (u) = \lambda\) ;
\item
  Let \(u \in \mathcal{D}_{\mathrm{B}}(\mathrm{E})\) and let \(\lambda\) be the family of its eigenvalues relative to B. The spectrum of \(u\) is the closure in K of the set of values of \(\lambda\) ;
\item
  If K = C, then the continuous functional calculus map of \(\mathcal{C}(\mathrm{Sp}(u))\) in \(\mathcal{L}(\mathrm{E})\) associates with a continuous function \(f \in \mathcal{C}(\mathrm{Sp}(u))\) the endomorphism \(\alpha(f \circ \lambda)\) in \(\mathcal{D}_{\mathrm{B}}(\mathrm{E})\) .
\end{enumerate}

By Lemma 1, the family of eigenvalues of an endomorphism diagonal in B is bounded, and the map from \(\mathcal{D}_{\mathrm{B}}(\mathrm{E})\) into \(\mathcal{B}(\mathrm{I})\) thus defined is a continuous isometric morphism of unital Banach algebras.

Let \(i \in I\) . For every \(j \in I\) , we have \(\langle u^{*}(e_{i}) | e_{j} \rangle = \lambda_{j} \langle e_{i} | e_{j} \rangle = \langle \overline{\lambda_{j}} e_{i} | e_{j} \rangle\) . It follows that \(u^{*}(e_{i}) = \overline{\lambda}_{i} e_{i}\) . The adjoint of u is therefore the diagonal endomorphism in the basis B for which \((\overline{\lambda}_{i})_{i \in I}\) is the family of eigenvalues. Assertion a) follows from this.

Let \(\lambda = (\lambda_i)_{i\in \mathrm{I}}\in \mathcal{B}(\mathrm{I})\) . For every \(x\in \mathrm{E}\) the family \((|\langle e_i|x\rangle |^2)_{i\in \mathrm{I}}\) is summable, with sum \(\| x\|^2\) (EVT, V, p. 22, prop. 5), whence, for every finite subset J of I:

\[
\begin{array}{r l r} & & {\left\| \sum_ {i \in \mathrm{J}} \lambda_ {i} p _ {i} (x) \right\| ^ {2} = \sum_ {i \in \mathrm{J}} | \lambda_ {i} | ^ {2} | \langle e _ {i} | x \rangle | ^ {2} \leqslant \Bigl (\underset {i \in \mathrm{I}} {\sup} | \lambda_ {i} | ^ {2} \Bigr) \sum_ {i \in \mathrm{J}} | \langle e _ {i} | x \rangle | ^ {2}} \\ & & {\leqslant \Bigl (\underset {i \in \mathrm{I}} {\sup} | \lambda_ {i} | \Bigr) ^ {2} \| x \| ^ {2}.} \end{array}
\]

Consequently, the family \((\lambda_{i}p_{i})\) is summable in \(\mathcal{L}(\mathrm{E})\) endowed with the topology of simple convergence. Its sum \(u_{\lambda}\) satisfies \(u_{\lambda}(e_{i}) = \lambda_{i}e_{i}\) ; it is therefore a diagonal endomorphism in the basis B, with eigenvalues \(\lambda\) . Assertion b) follows from this.

The last assertions follow from \(a\) ), from example 3 of I, p. 17 and from example 4 of I, p. 111.

Remark. --- The Banach algebra \(\mathcal{D}_{\mathrm{B}}(\mathrm{E})\) is a maximal commutative closed subalgebra of \(\mathcal{L}(\mathrm{E})\) . Indeed, let u be an endomorphism of E that commutes with \(\mathcal{D}_{\mathrm{B}}(\mathrm{E})\) . Let \(i \in I\) . Since the orthogonal projector \(p_{i}\) is diagonal in the basis B, we have \(p_{i}(u(e_{i})) = u(p_{i}(e_{i})) = u(e_{i})\) , which implies that \(u(e_{i})\) is proportional to \(e_{i}\) . Thus u is diagonal in the basis B.

If E is infinite-dimensional, there exist maximal commutative involutive subalgebras of \(\mathcal{L}(\mathrm{E})\) that are not isomorphic to \(\mathcal{D}_{\mathrm{B}}(\mathrm{E})\) (exercise 5 of IV, p. 314).

PROPOSITION 2. --- Let \(u \in \mathcal{D}_{\mathrm{B}}(\mathrm{E})\) and \(\lambda = (\lambda_i)_{i \in \mathrm{I}}\) be the family of its eigenvalues.

\begin{enumerate}
\def\labelenumi{\alph{enumi})}
\tightlist
\item
  The following conditions are equivalent:
\end{enumerate}

\begin{enumerate}
\def\labelenumi{(\roman{enumi})}
\item
  One has \(\lambda \in \mathcal{C}_{0}(\mathrm{I};\mathrm{K})\) ;
\item
  The endomorphism u is compact;
\item
  The family \((\lambda_{i}p_{i})_{i\in\mathrm{I}}\) is summable in the Banach space \(\mathcal{L}(\mathrm{E})\) . Its sum is then equal to u.
\end{enumerate}

\begin{enumerate}
\def\labelenumi{\alph{enumi})}
\setcounter{enumi}{1}
\tightlist
\item
  Suppose that \(u\) is compact and let \(\Lambda\) denote the set of values of \(\lambda\) . The set of \(i \in I\) such that \(\lambda_i \neq 0\) is countable. If E is infinite-dimensional, we have \(\mathrm{Sp}_{\mathrm{s}}(u) = \Lambda - \{0\}\) and \(\mathrm{Sp}(u) = \Lambda \cup \{0\}\) . If E is finite-dimensional, then \(\mathrm{Sp}_{\mathrm{s}}(u) = \mathrm{Sp}(u) = \Lambda\) .
\end{enumerate}

By Lemma 1, we have \(\| \sum_{j\in J}\lambda_jp_j\| = \sup_{j\in J}|\lambda_j|\) for every finite subset J of I. It follows, by the Cauchy criterion, that the family \((\lambda_ip_i)\)

is summable in \(\mathcal{L}(\mathrm{E})\) if and only if the family \(\lambda\) tends to 0 at infinity, which implies that conditions (i) and (iii) are equivalent.

Condition (iii) implies that \(u\) is compact (corollary of proposition 2 of III, p. 4). Conversely, suppose that the endomorphism \(u\) is compact. Since the family B is bounded in E, its image under \(u\) in E is relatively compact in E (III, p. 2), hence precompact. Let \(\varepsilon > 0\) and let \(J\) be a finite subset of I such that the image of B under \(u\) is contained in the union of balls of radius \(\varepsilon\) and centers \(u(e_j)\) for \(j \in \mathrm{J}\) . Let \(i \in \mathrm{I} - \mathrm{J}\) . There exists \(j \in \mathrm{J}\) such that \(\| u(e_i) - u(e_j) \| \leqslant \varepsilon\) , hence

\[
\left| \lambda_ {i} \right| ^ {2} \leqslant \left| \lambda_ {i} \right| ^ {2} + \left| \lambda_ {j} \right| ^ {2} = \left\| u \left(e _ {i}\right) - u \left(e _ {j}\right) \right\| ^ {2} \leqslant \varepsilon^ {2}.
\]

Consequently, we have \(\lambda \in \mathcal{C}_0(\mathrm{I};\mathrm{K})\) , that is, condition (i).

Finally, the spectrum of u and its sensitive spectrum are computed in terms of \(\lambda\) using prop. 1, c) and proposition 5 of III, p. 90.

Remark. --- When I is infinite, condition \(\lambda\in\mathcal{C}_{0}(\mathrm{I};\mathrm{K})\) can also be stated as “the family \(\lambda\) tends to 0 along the filter of complements of finite subsets of I.”

